\section{Conclusion}
ZeroCausal presents a shift in provenance-based IDS design by removing the clean-data and label assumptions. By leveraging online causal discovery via Tigramite PCMCI and SCM residual scoring, ZeroCausal detects stealthy APT attacks purely as violations of learned causal mechanisms. Combined with online conformal prediction, dynamic change-point-triggered SCM refitting, and rolling calibration queue updates, ZeroCausal achieves an average AUC of \textbf{0.727 $\pm$ 0.071} on real OpTC logs over 10 seeds and \textbf{0.9656} on BETH (cross-platform Linux host telemetry), with conformal false-alarm bounds tightly controlled at 4.09\% $\pm$ 0.66\% and graceful concept-drift adaptation. However, operational evaluation at the target 5\% FPR budget highlights a key challenge: ZeroCausal achieves a recall of 24.21\% (F1: 0.3286), missing stealthy attacks that blend into baseline variance. Our evaluation across five diverse benchmarks---including an honest failure case on StreamSpot (AUC 0.4991) that diagnoses limitations of linear SCMs on heterogeneous graph-type provenance data---demonstrates both the strengths and current boundaries of causal invariance-based detection. Hardening the calibration step against contamination and scaling to high-dimensional streams remain the primary directions for future work.
